\section{Published POAF Evidence and Why High Scores Mislead}
\label{sec:literature}

A central obstacle in interpreting cardiac-surgery machine learning papers is that published POAF discrimination spans from 0.58 to 0.98~\cite{fields2025poaf,lopeztuesta2026poaf}. Without separating endpoint definitions, prediction windows, surgical case mix, predictor modalities, and validation designs, a 0.68 pre-incision national-registry benchmark can appear deceptively weak beside single-center 0.85--0.97 reports. We searched PubMed (566 unique records through November 2026), the STS publication index, and backward citations from recent systematic reviews (full 42-study annotated inventory in Appendix~A).

\subsection{What STS Cohorts and Systematic Reviews Establish}
The 2018 STS adult cardiac surgery risk models focus on operative mortality and major morbidity~\cite{shahian2018sts,obrien2018sts2}, while related STS-based CABG machine-learning models target operative mortality, major morbidity, cost, and readmission~\cite{zeavera2022stsml}. In the closest multi-center STS isolated-CABG evaluation of bedside POAF scores (8,976 cases across five centers), Pollock et al.~\cite{pollock2018sts} reported AUCs of \textbf{0.68} for CHARGE-AF~\cite{alonso2013chargeaf}, \textbf{0.66} for age alone, \textbf{0.65} for the Mariscalco POAF score~\cite{mariscalco2014poaf}, \textbf{0.59} for $\mathrm{CHA_2DS_2\text{-}VASc}$~\cite{lip2010cha2ds2vasc}, and \textbf{0.58} when the STS predicted mortality score was repurposed as a POAF comparator---and even those regressions fitted a pooled random-intercept model across the five centers rather than testing strictly unseen hospitals. Modest pre-incision discrimination recurs across independent registries: Segar et al.~\cite{segar2023sts} observed a best bedside score AUC of 0.65 under STS criteria, Potdar et al.~\cite{potdar2022offpump} reported 0.64 for the STS risk score in prospective off-pump CABG ($N=1{,}108$), Thor{\'e}n et al.~\cite{thoren2012cabg} and El-Chami et al.~\cite{elchami2012poaf} reported 0.60--0.62 in large CABG series, Chung et al.~\cite{chung2023uk} reached 0.68 internally across 33,464 UK cardiac operations, Rosati et al.~\cite{rosati2025polaris} reported 0.635 for the two-center POLARIS score, and prospective external validation by Cameron et al.~\cite{cameron2018validation} ($N=1{,}416$ CABG/valve patients) yielded external AUCs of 0.651 (POAF score), 0.593 ($\mathrm{CHA_2DS_2\text{-}VASc}$), and 0.563 (AFRisk) across eligible subsets.

Recent systematic reviews confirm how sharply discrimination shrinks outside derivation cohorts. Across cardiac-surgery POAF models appraised under PROBAST~\cite{wolff2019probast}, Fields et al.~\cite{fields2025poaf} reported a \textbf{median apparent AUC of 0.71} (range 0.60--0.98) falling to a \textbf{median external AUC of 0.61} (0.51--0.77), rating every evaluated analysis at high risk of bias. L{\'o}pez Tuesta et al.~\cite{lopeztuesta2026poaf} reviewed 39 studies, finding only six external validations, inconsistent calibration reporting, and 32 of 39 at high risk of bias. Pandey et al.~\cite{pandey2023poaf} pooled AUCs of 0.60 (STS score), 0.66 (POAF score), and 0.68 ($\mathrm{CHA_2DS_2\text{-}VASc}$). Conversely, Eini et al.~\cite{eini2026cabg} reported a pooled CABG model AUC of 0.84 with extreme heterogeneity (generalized $I^2 = 87.3\%$) across heterogeneous cohorts, predictors, and validation designs.

\begin{table*}[t]
\centering
\tablefont
\caption{\ct{Published POAF prediction benchmarks and why higher reported AUCs are not directly comparable to pre-incision isolated-CABG unseen-hospital STS transport.} Extended 42-study annotated inventory appears in Appendix~A.}
\label{tab:lit_landscape}
\setlength{\tabcolsep}{4.3pt}
\begin{tabularx}{\textwidth}{@{}L{3.25cm}L{2.45cm}L{1.8cm}R{1.45cm}Y@{}}
\toprule
\multicolumn{1}{@{}l}{\thd{Study / Review}} & \multicolumn{1}{l}{\thd{Cohort \& Surgery}} & \multicolumn{1}{l}{\thd{Window}} & \multicolumn{1}{r}{\thd{Reported AUC}} & \multicolumn{1}{l@{}}{\thd{Key Methodological Factor Governing Comparison}} \\
\midrule
\multicolumn{5}{@{}l@{}}{\textit{\color{sub}Closest STS / Multi-Center Pre-Incision Baselines \& Systematic Reviews:}} \\
Pollock et al.\ (2018)~\cite{pollock2018sts} & STS, 5 centers, iso-CABG ($N=8{,}976$) & Pre-incision & 0.58--0.68 & CHARGE-AF~\cite{alonso2013chargeaf} 0.68, age 0.66, POAF 0.65, STS mort 0.58; pooled site random-effect fit (not unseen sites). \\
Mariscalco et al.\ (2014)~\cite{mariscalco2014poaf} & 3 EU centers, mixed ($N=17{,}262$) & Pre-incision & 0.64 (0.65 log.) & Bedside point score 0.64 in derivation/validation; includes valve indicator (23\% CABG vs.\ 40\% combined). \\
Cameron et al.\ (2018)~\cite{cameron2018validation} & Prospective ext.\ ($N=1{,}416$) & Pre-incision & 0.56--0.65 & External validation: POAF 0.651, $\mathrm{CHA_2DS_2\text{-}VASc}$ 0.593, AFRisk 0.563 in eligible subsets. \\
Fields et al.\ (2025)~\cite{fields2025poaf} & Systematic review (PROBAST~\cite{wolff2019probast}) & Mixed & 0.61 ext.\ (0.71 app.) & Median external AUC 0.61 (0.51--0.77) vs.\ 0.71 apparent; 100\% high risk of bias. \\
L{\'o}pez Tuesta (2026)~\cite{lopeztuesta2026poaf} & Systematic review (39 studies) & Mixed & 0.60--0.98 & Only 6/39 externally validated; 32/39 rated high risk of bias. \\
\midrule
\multicolumn{5}{@{}l@{}}{\textit{\color{sub}Representative Higher-AUC Reports Driven by Later Windows, Case Mix, Small-$N$ Optimism, or Richer Modalities:}} \\
Magee et al.\ (2007)~\cite{magee2007atrial} & STS iso-CABG ($N=19{,}083$) & Intra/post-op & 0.72 & Uses cardiopulmonary bypass and $\ge$24\,h postoperative ventilation (unavailable by incision). \\
Mathew et al.\ (2004)~\cite{mathew2004multicenter} & 70 centers, CABG+valve ($N=4{,}657$) & Intra/post-op & 0.77 & Incident-POAF model ($n=1{,}564$ val.) includes concomitant valve surgery and post-op drug withdrawal (cf.~\cite{waldron2018poaf}). \\
Cho et al.\ (2026)~\cite{cho2026poaf} & Mixed cardiac ($516 \to 1{,}701$) & Perioperative & 0.891 $\to$ 0.726 & Sharp internal-to-external drop illustrates site transport difficulty; mixed surgery. \\
Ma et al.\ (2025)~\cite{ma2025noaf} & CABG, 2 sites ($N=508$ ext.) & Post-procedure & 0.931 ext. & Legitimate external test, but requires actual surgical duration and LA diameter. \\
Javaherforooshzadeh et al.\ (2026)~\cite{javaherforooshzadeh2026car} & Single center CABG ($N=109$) & Pre-op lab & 0.971 app. & Apparent discrimination on 29 events; no external validation reported. \\
Han et al.\ (2025)~\cite{han2025aiecg} & Cardiac surgery + 12-lead ECG & Pre-op AI-ECG & 0.643 $\to$ 0.680 & Waveform AI-ECG adds signal beyond clinical POAF score (0.901 is non-surgical AF task). \\
Khalpey \& King (2026)~\cite{khalpey2026digital} & STS + EHR multimodal ($N=38{,}742$) & Multimodal & 0.89 & Combines continuous ECG, hemodynamics, and biomarkers absent from STS structured fields. \\
\midrule
\textbf{\color{accent}This study (EvoTarget-FT)} & \textbf{STS iso-CABG (74 unseen sites, $N=21{,}241$)} & \textbf{By incision (207 raw)} & \textbf{\color{accent}0.6763} & \textbf{Strictly pre-incision STS structured predictors; 0.0\% hospital overlap; 74-hospital cluster CIs.} \\
\bottomrule
\end{tabularx}
\end{table*}


\subsection{Five Structural Reasons Published High AUCs Are Non-Comparable}
Table~\ref{tab:lit_landscape} classifies representative high-AUC POAF reports by the exact factor preventing direct comparison to pre-incision isolated-CABG transport across unseen STS centers:
\begin{enumerate}
  \item \textbf{Intraoperative and postoperative predictor availability:} The two most cited classical algorithms---Magee et al.'s STS model (AUC 0.72, $N=19{,}083$)~\cite{magee2007atrial} and Mathew et al.'s multicenter incident-POAF index (AUC 0.77 in 1,564 validation patients from 4,657 CABG/valve cases)~\cite{mathew2004multicenter}---incorporate cardiopulmonary bypass use, prolonged ventilation ($\ge 24$\,h), concomitant valve surgery, or postoperative medication administration and withdrawal. Waldron et al.~\cite{waldron2018poaf} observed higher discrimination for McSPI (0.77), which incorporates later perioperative data, than for strictly preoperative scores (0.58--0.66) evaluated in the same cohort. Recent ML models achieving 0.931 externally (Ma et al.~\cite{ma2025noaf}, using actual surgical duration), 0.813 after the first 24\,h of ICU data (Cao et al.~\cite{cao2025mimic4}), relying on cross-clamp time (Parise et al.~\cite{parise2023cabg}, whose 0.95/0.94 AUCs are resampling maxima vs.\ 0.79 held-out accuracy), or sampling postoperative pericardial fluid (Luo et al.~\cite{luo2026pcmpo}, 0.865) address post-procedure surveillance rather than pre-incision risk.
  \item \textbf{Mixed valve and CABG case mix and setting:} POAF incidence rises sharply by procedure class---23\% in isolated CABG, 31\% in isolated valve, and 40\% in combined procedures in Mariscalco et al.~\cite{mariscalco2014poaf} (approximately 20\%--30\%, 30\%--40\%, and 40\%--50\% in the 2026 STS guideline~\cite{chatterjee2026sts}; 24\% pooled CABG incidence across 207 studies~\cite{farzaneh2026incidence}). Because surgical class itself strongly separates risk, mixed cardiac cohorts can change achievable discrimination (e.g., Burgos COM-AF~\cite{burgos2021comaf}, Cho et al.~\cite{cho2026poaf}, and Karri et al.'s preoperative MIMIC-III ICU model at 0.74~\cite{karri2021poaf}).
  \item \textbf{Single-center optimism and small-$N$ instability:} Cho et al.~\cite{cho2026poaf} dropped from 0.891 internally ($n=516$) to 0.726 externally ($n=1{,}701$); Shakhgeldyan et al.~\cite{shakhgeldyan2026cabg} dropped from 0.999 training to 0.76--0.78 on a single-center test; Peng et al.~\cite{peng2024codaaf} reported 0.894 on a single-center internal CODA-AF split using CRP and left atrial diameter; and a 109-patient CRP/albumin study reported 0.971 apparent AUC with 29 events and no external validation~\cite{javaherforooshzadeh2026car}.
  \item \textbf{Richer non-registry modalities or composite endpoints:} Waveform 12-lead AI-ECG lifts the POAF score from 0.643 to 0.680 in a separate surgical cohort (5,204 ECGs, 2,266 patients, 30-day ECG endpoint; Han et al.~\cite{han2025aiecg}, where 0.901 belongs to the non-surgical AF detection task) and reaches 0.66 standalone in SURGICAL-AF~\cite{gruwez2024aiecg}; genomic panels add signal over clinical baselines (Kolek et al.~\cite{kolek2015genetic}, $0.69 \to 0.72$; Hua et al.~\cite{hua2025genomic}, 0.81 internal $\to$ 0.76 independent validation with bypass covariates). Multimodal digital-twin work by Khalpey and King~\cite{khalpey2026digital} reached 0.89 (vs.\ 0.72 for their STS score comparator across 38,742 mixed cardiac operations) by combining continuous ECG, hemodynamics, and biomarkers across linked EHR and registry sources, while Shen et al.~\cite{shen2025surgvae} reported 0.831 as a six-outcome macro-AUROC rather than POAF alone.
  \item \textbf{Endpoint ascertainment:} STS captures treatment-requiring AF/flutter; telemetry-defined or 30-day ECG endpoints alter both incidence and predictability~\cite{filardo2017telemetry,filardo2018onset,filardo2020telemetry,perezgrovas2023jaha}.
\end{enumerate}
We therefore interpret absolute AUCs strictly in the context of prediction time, surgical eligibility, outcome ascertainment, and validation design---which immediately motivates auditing our national STS cohort for schema leakage and hospital overlap (\S\ref{sec:cohort_splits}).
